\section{A visible label need not identify the observation mechanism}
\label{sim:bound-observation-map}

In the preceding example, a first message $H$ requires pay under one interpretation and draw under another. An additional label can resolve the conflict if it tells the actor which interpretation applies. Suppose, however, that both interpretations attach the same label, \emph{reviewed}, to that message. The actor sees a label but still cannot distinguish the two cases. Declaring a label visible and declaring a mechanism identified are different assumptions.

The distinction changes the feasible policy set. With distinct labels supplied before the first decision, the two four-policy problems admit 16 joint policies. With the same label supplied in both, the conflicting instructions remain and no common policy is feasible. A simulator that automatically inserts its internal mechanism identifier into the actor's message would report the first answer while pretending to solve the second problem.

The selected observation declaration therefore records the actual token presented to the actor. It permits distinct mechanisms to emit the same token. Source identities and temporal proof labels remain attached to the analyst's explanation, but do not enter the decision history unless the declaration supplies a separate observation with that meaning. The numerical comparison keeps these cases apart even though both could informally be described as a labeled source.

\subsection{A message map must use information available at that date}

Hiding internal identifiers is insufficient if the message itself encodes an unavailable fact. A map that takes the entire financial history as its input could send one message when a hidden borrowing right exists and another when it does not. The resulting policy would exploit that right even if the source observations never revealed it. Giving the message an innocuous name would not remove the additional information.

The finite reference avoids that ambiguity by restricting what each map can read. At the first date, the coarse view supplies the market signal; the refined view supplies that signal and the current coupon-trigger observation. At the second date, the map receives the already specified trigger and cash-coverage observations. The actor retains the first message and its own earlier instruction. Hidden contractual rights, complete world identifiers and later observations cannot be introduced as new first-date map inputs.

For each selected source outcome, the declaration must map every supported observation key to a message. A missing key is an unfinished declaration, not permission to remove the corresponding financial history. Two entries for the same key are also invalid. These restrictions establish a closed finite interpretation of the existing observations. They do not establish that those observations describe a particular company's actual information practices, nor do they define a general language for arbitrary messages.

This ordering is useful when designing richer models. First specify which quantities can be observed and when; then specify how the observation is represented to the actor; finally constrain decisions to agree on indistinguishable histories. Reversing the order can make a desirable policy feasible by designing messages around the hidden state that policy would need to know.

\subsection{Source, actor and time remain separate premises}

A claimed visible tag carries an actor, the selected source revision, its content identity, a token and an observation time. The first decision must agree exactly with the temporal family's cutoff; the second occurs strictly later. A tag claimed at the first decision must have been available by that instant. Equivalent fixed-offset spellings describe the same instant, whereas a tag received one ten-millionth of a second afterward is late under the selected exact-time convention.

Unknown tag availability leaves the declaration unresolved. So does a supplied observation whose chronology cannot be established from the retained authorship bounds. The current reference treats this conservatively: it does not discard inconvenient clock possibilities or silently refine the source family with a new tag observation. A later method may condition or refine that family under explicitly admitted additional evidence, but it must preserve the distinction between the original question and the revised one.

These unresolved cases differ from a complete observation model with no feasible policy. In the latter, the available histories are specified and every possible instruction fails somewhere in the retained financial family. In the former, the model has not established which histories the actor may distinguish. Returning the same empty policy list for both would hide a missing information premise behind an apparently decisive financial conclusion.

A content identity detects alteration relative to expected bytes; it does not show that the actor received them. Likewise, the synthetic tag records used in the comparison test the consequences of supplied assumptions. They are not observations of a real company's managers or investors. That distinction remains necessary even when every mapping and timestamp check passes.

\subsection{Reuse the calculation only when the decision problem is unchanged}

A corrected source can preserve the same financial meaning and messages. The numerical policy calculation can then be reused while its source explanation changes. A different actor can also have the same supplied observation map; the actor binding must change even if the resulting mathematical problem is identical. Neither case licenses removal of the old and new source supports from the retained temporal family.

Changing a message label is different. A policy is a rule for responding to the labels the actor receives. Relabeling $H$ as another token changes that rule's interface even if four policies remain feasible and a mathematician can construct a one-to-one relabeling between the sets. The selected reference recomputes this case. It does not claim to have proved and implemented a general translation between equivalent policy representations.

The reuse comparison consequently includes the complete financial and action population, the actual first and second messages, the meanings of the actions, and the numerical methods. Identical repeated constraints can share computation while keeping every source outcome and temporal proof support attached. Distinct problems with zero feasible policies retain distinct identities. A future correction that is unavailable at the earlier cutoff remains in the source record without replacing the earlier view.

The independent checks reconstruct the messages from the supplied declarations and original observations, compare full policy sets from two formulations, and reject missing or duplicated keys, hidden-state inputs, wrong actors or times, altered sources and forged prior results. They also check changed second-date labels without changing first-date messages. These examples make the dependency rule concrete: identical counts are not enough, but a changed source explanation need not force an unchanged numerical problem to be solved again.

For this observation-binding profile, the remaining computational difficulty is how to continue an unfinished temporal search while retaining its unexamined possibilities. Section~\ref{sim:release-temporal-frontier} develops that extension for the selected temporal domain. The conservative full-domain remainder already prevents a partial result from masquerading as complete; a checked, resumable partition would additionally preserve useful completed work. Actual message provenance, richer observation languages and refinement by new tag evidence remain distinct extensions.\footnote{The finite declarations, complete mapping comparisons, 22 main cases and 30 adversarial controls are retained in \path{data/accounting-master/observation-map-binding-20261027/}. A separate saved audit rechecks all main cases and four additional source, exact-time and token cases. The reference retains full source supports under numerical sharing and makes no claim of actual actor knowledge or authenticated evidence.}
